% ======================================================================
%  Section 4.9 --- Future work.  WRITTEN.
%  This is the sec:future that 4.4.4, 4.7.1, 4.7.2, 4.7.3 and 4.8 all
%  refer to. Five dangling references resolve when you paste it.
%
%  WHERE TO PUT IT. All five references are from inside Chapter 4 and all
%  five are about Chapter 4's own gaps, so it belongs at the end of
%  Chapter 4, after the conclusion. If your thesis gathers all future work
%  into the general conclusion instead, move the whole section there and
%  KEEP THE LABEL sec:future -- the five references will still resolve,
%  and only the printed section number changes.
%
%  HOW IT IS ORDERED. By value, not by topic, and the first item is worth
%  more than the other eight together: the chapter measured a three-order-
%  of-magnitude accuracy gain from a gain schedule that does not vanish,
%  which is a larger effect than anything else in it. An examiner who
%  reads only the first paragraph should come away with that.
%
%  ONE THING I DID NOT LIST. Monitoring the conditioning of H_D. Section
%  4.5, as corrected, argues that the floor makes the rank condition
%  non-binding by construction, so the measurement is not needed; listing
%  it here as outstanding work would undercut that argument. I mention it
%  only as an option, and only in the parenthesis where it belongs.
%
%  TWO ITEMS ARE MINE RATHER THAN YOURS, and you should decide whether you
%  want them:
%    - the reference-structure comparison in "Distinguishing occlusion
%      from absent texture", which is a concrete research idea and not
%      merely a gap to be filled. It follows directly from 4.4.4's
%      acknowledged limitation and it is cheap. If you would rather the
%      thesis not propose specific solutions, cut the second half of that
%      paragraph and keep the statement of the problem.
%    - the correlated-residuals paragraph, which names a theoretical
%      misspecification in the estimator rather than an experiment. It is
%      the kind of thing a sharp examiner raises, and having raised it
%      yourself is better than being asked.
% ======================================================================

\section{Future work}
\label{sec:future}

Three kinds of work remain outstanding: measurements that the claims of this
chapter require but do not yet have, an improvement that the experiments
identified inadvertently, and extensions of the method itself. They are given
below in order of the return expected from each.

\paragraph{A gain schedule that does not vanish with the error.}
This is the most valuable improvement the chapter identifies, and the evidence
for it is already in hand. The adaptive gain~\eqref{eq:lambda} is proportional
to the intensity error, and the commanded velocity is proportional to the
residual in addition, so the velocity decays roughly as the cube of the error
and the nominal runs stall while a sub-millimetre misalignment remains. The
occluded runs, in which a fixed residual persists and the gain is consequently
held some fifty-four times higher, reach the desired pose three orders of
magnitude more closely. That comparison is confounded, in that the two
configurations differ in more than the gain, but the mechanism is established
independently: the gain and the damping are both driven by the intensity
error, and the fifth-power relation between them is satisfied to within
$0.62\%$ across all three runs. What is needed is a direct test. Three
schedules are worth comparing on the nominal configuration, where the stall
occurs: a floor on $\lambda$; a gain normalised by the current residual rather
than proportional to it, so that the step length rather than the velocity is
controlled; and a two-phase schedule that switches from the present law to a
constant gain once the damping has collapsed. Each is a few lines of the
implementation, and the chapter's own measurement suggests the return may be
of the order of three orders of magnitude in final accuracy. It should be
noted that decoupling the gain from the error in this way removes the
mechanism by which the present scheme terminates, so a stopping criterion
would have to be supplied in its place --- which the chapter should have had
in any case, the inherited absolute threshold on $\norm{\mathbf{e}}^{2}$ never
having been met in any run.

\paragraph{A study over several initial poses.}
This is the largest evidential gap. Every figure reported in
Section~\ref{sec:ablation-results} rests on a single run from a single initial
displacement, which is why no ordering of the three variants is asserted
anywhere in the chapter. The two weighted variants differ by $28\%$ in
translation in favour of the residual-based weighting and by $8\%$ in rotation
in favour of the structure-informed one, and neither margin can be shown to be
systematic on one run each. A study over a sample of initial poses --- varying
the direction and magnitude of the displacement, and in particular the depth
component, since Section~\ref{sec:nominal} finds the degradation concentrated
there --- would allow the comparison to be reported as a distribution rather
than as a pair of numbers, and would either establish the advantage claimed
for the structural information or refute it. Nothing else in this list would
change the standing of the chapter's central question as much.

\paragraph{Measurements outstanding.}
Three are cheap and each closes a specific gap. The first is a weight map for
Variant~2 under the same conditions as those of
Section~\ref{sec:weight-maps}: the chapter can state that the structure
normalisation does not eliminate the early rejection of valid measurements,
$13.6\%$ of them being held below a quarter weight at the first iteration, but
it cannot state whether the normalisation reduces that figure relative to a
purely residual-based estimator, no map having been recorded for the latter.
One run supplies the comparison. The second is a proper measurement of the
computational overhead of the weighting, the timings in the present logs being
dominated by rendering and machine load to the point of reversing the
algorithmic ordering; a clock placed around the weight computation alone, on
an unloaded machine, would test the prediction of
Section~\ref{sec:hermite-principle} that the six additional convolutions cost
some twenty per cent over the twenty-eight the feature already requires. The
third is the Frobenius norm of the iteration-to-iteration change in the
weighting, $\lVert\mathbf{D}_{k}-\mathbf{D}_{k-1}\rVert_{F}$, logged against
the iteration. Section~\ref{sec:robust-stability} assumes the weights vary
slowly relative to the control and does not verify it; the terminal jitter
reported in Section~\ref{sec:occlusion}, in which the structure-informed
variant retains a commanded velocity thirty-four times that of the
residual-based one after its pose has ceased to improve, suggests the
assumption deserves checking. (The conditioning of
$\mathbf{H}_{\mathbf{D}}$ could be logged at the same time, though
Section~\ref{sec:robust-stability} argues that the floor
of~\eqref{eq:hessian-floor} makes the rank condition non-binding by
construction and the measurement therefore unnecessary.)

\paragraph{Characterising the admissible occlusion fraction.}
The chapter tests one occluder, of one shape, at one position, covering
$5.80\%$ of the measurements at the desired pose. The robust framework of
Section~\ref{sec:robust-framework} gives an M-estimator a finite breakdown
point, so there exists a fraction of corrupted measurements beyond which each
variant fails, and locating it is a well-posed experiment: sweep the occluder
area from zero upward and record the fraction at which the final pose error
departs from its unoccluded value. The result would be of direct practical
use, since it is the quantity a practitioner needs in order to decide whether
a given scene is within the method's competence, and it would test whether the
structural information extends the admissible fraction --- which is a
different and more tractable question than whether it improves the accuracy at
a fixed fraction, and one on which the two weightings might well be
distinguishable where the present experiment cannot distinguish them.

\paragraph{Distinguishing occlusion from absent texture.}
The structure measure reports that a region is locally flat, not why, and
Section~\ref{sec:hermite-weight-limitation} accepts that a genuinely
untextured but visible region is suppressed along with an occluded one. The
reference image offers a way to separate the two cases that the present scheme
does not use. The structure measure can be evaluated on the reference as well
as on the current image, and the two compared at corresponding positions: a
region flat in both is scene texture that is genuinely absent, whereas a
region flat in the current image and structured in the reference is occluded.
The comparison requires one additional evaluation of the measure, performed
once on the reference and reused throughout, so the cost is negligible. Its
limitation should be stated plainly: the two images are in spatial
correspondence only as the pose converges, so the criterion is reliable in
just the regime in which the present mechanism already works, and it would
sharpen the discrimination rather than extend it to the early iterations. A
textured occluder would defeat it as it defeats the present scheme, since such
an occluder produces a large response in the current image and would be
retained.

\paragraph{Reducing the early over-rejection.}
Section~\ref{sec:tukey-limitation} identifies the rejection of valid
measurements while the pose error is large as the characteristic defect of a
structure-blind estimator, and Section~\ref{sec:weight-maps} measures it at
$13.6\%$ of the clean scene at the first iteration with the structure
normalisation active. The normalisation redistributes the residuals relative
to one another but cannot alter the fact that, far from the solution, the
residual is large almost everywhere and the distribution is nothing like the
one the estimator presumes. Two remedies are standard elsewhere and untried
here. The tuning constant of the Tukey function can be annealed, beginning
large so that almost nothing is rejected and contracting to its nominal value
of $4.6851$ as the error falls, which is the deterministic-annealing strategy
used in robust registration; alternatively the weighting can simply be held at
the identity for an initial phase and enabled once the residual distribution
has narrowed. Either would be straightforward to evaluate against the weight
maps already produced, which give the figure to be improved upon.

\paragraph{Real images and other classes of outlier.}
Every result reported here is obtained in simulation, with the target rendered
from a photograph by a planar image simulator, so the only departures from the
constancy assumption are those deliberately introduced. Experiments on a real
camera would add sensor noise, calibration error and, most importantly,
illumination that varies over the scene and over time. That last is not merely
a further difficulty but the test of a specific claim made in
Section~\ref{sec:hermite-weight-limitation}: the intensity-based structure
measure was rejected there in favour of the zero-mean formulation although it
attained the smaller error on the nominal run, $0.088^{\circ}$ against
$0.385^{\circ}$ in rotation, on the ground that it would behave differently on
a scene in which brightness and texture were uncorrelated. Simulation from a
single photograph cannot produce such a scene; a real one will. In the same
direction, specular highlights, moving objects and partial illumination change
are all local violations of the same assumption as occlusion but with
different spatial and temporal signatures, and there is no reason to expect a
weighting tuned against a static occluder to handle all of them equally.

\paragraph{Correlated residuals.}
A theoretical point deserves recording, since it bears on the estimator rather
than on the experiments. The measurements are Hermite responses computed with
$9\times9$ filters at every interior position, so the supports of neighbouring
measurements overlap heavily and their residuals are strongly correlated in
space. The M-estimator of Section~\ref{sec:robust-framework}, like the robust
scale of Section~\ref{sec:mad}, treats them as independent: the median
absolute deviation of $264\,000$ correlated residuals is not the estimate of
scale that the same expression would give for independent ones, and the
effective number of independent measurements is smaller, perhaps considerably
smaller, than $P$. This does not invalidate the scheme, which works, but it
does mean that the breakdown point quoted for the estimator is an
overstatement and that the weighting is not the maximum-likelihood
construction that the derivation suggests. Quantifying the effective
independence --- for instance from the autocorrelation of the residual field
--- would put the robust scale on a sounder footing and would give the
occlusion fraction of the preceding paragraph a theoretical prediction to be
tested against.

\paragraph{Learned measures of reliability.}
Finally, the structure measure adopted here is hand-designed, and its
justification is that it approximates a quantity --- the norm of a row of the
interaction matrix --- that can be written down. An alternative is to learn a
per-measurement reliability directly, from pairs of images in which the
corrupted regions are known, and the appeal is that a learned measure could
condition on evidence the present one discards: the residual's history over
the trajectory, the response across scales rather than its aggregate, the
agreement of a measurement with its neighbours. The reservation is equally
plain, and this chapter is the reason for it. A measure with no analytical
relation to the interaction matrix cannot be defended by the argument of
Section~\ref{sec:hermite-principle}, and the ablation of
Section~\ref{sec:hermite-weight-limitation} shows what happens to a measure
that lacks such a relation: it can attain the better final error on the
configuration it was chosen on while resting on no property that transfers.
Any learned weighting should therefore be evaluated on the criteria this
chapter applied to a hand-designed one --- whether it recovers the identity as
the data cease to contain outliers, whether its behaviour depends on
brightness or on structure, and whether it survives a scene it was not
selected on.

% ======================================================================
%  EVERY NUMBER HERE IS FROM THE CHAPTER. No new quantity is introduced.
%    54, 0.62 per cent          gain and damping      (4.7.2)
%    three orders of magnitude  nominal vs occluded   (4.7.2)
%    28 per cent, 8 per cent    V2/V3 margins         (4.7.2)
%    thirty-four                terminal jitter       (4.7.2)
%    13.6 per cent              early over-rejection  (4.7.3)
%    5.80 per cent              occluder at desired pose (4.6.3, 4.7.2)
%    twenty per cent, six, twenty-eight  convolution cost (4.4.1)
%    0.088 vs 0.385 deg         measure ablation      (4.4.4)
%    4.6851                     Tukey constant        (4.3.1)
%    264 000                    feature dimension     (4.6.2)
%
%  REFERENCES THIS SECTION RESOLVES
%    4.4.4  "...is identified as an open point in Section~\ref{sec:future}"
%    4.7.1  "...the study over several initial poses that
%            Section~\ref{sec:future} identifies as outstanding"
%    4.7.2  "...identified in Section~\ref{sec:future}" (twice: the
%            multi-pose study, and the gain schedule)
%    4.7.3  via 4.4.4
%    4.8    "...identified in Section~\ref{sec:future}"
%  Check each of these reads correctly against the paragraph it now points
%  to. The two in 4.7.2 point at different paragraphs of this section, so
%  if you would rather they were specific, give the paragraphs their own
%  labels -- sec:future-gain and sec:future-poses would do.
%
%  LABELS USED ABOVE
%    eq:lambda , eq:hessian-floor
%    sec:robust-framework   4.2      sec:mad                4.3.2
%    sec:tukey-limitation   4.3.3    sec:hermite-principle   4.4.1
%    sec:hermite-weight-limitation 4.4.4
%    sec:robust-stability   4.5      sec:ablation-results    4.7
%    sec:nominal            4.7.1    sec:occlusion           4.7.2
%    sec:weight-maps        4.7.3
%
%  IF YOU NEED IT SHORTER, cut in this order: the learned-measures
%  paragraph (it is the most speculative), then the correlated-residuals
%  paragraph (it is the most technical), then the second half of the
%  occlusion-versus-texture paragraph. Do not cut the gain schedule or the
%  multi-pose study; the first is the chapter's most valuable finding and
%  the second is the answer to the question an examiner will certainly ask
%  about single runs.
% ======================================================================
